\section{Introduction}
\label{sec:intro}

Avenida L\'azaro C\'ardenas runs about ten kilometres from the Garza Sada
and Carretera Nacional interchange in south Monterrey to Avenida Gómez
Mor\'in in San Pedro Garza Garc\'ia. It carries traffic between the
residential south and the Valle Oriente business district, and its peaks
are long. The usual explanation goes like this: dozens of streets, ramps
and plaza entrances feed the avenue, and if you add up their lanes you get
well over a hundred, all emptying into five or six lanes per direction.

That argument is a lane count. A queueing system does not care how many
lanes arrive. It cares how many vehicles per hour arrive and how many per
hour can leave, at each point where they meet. A one-lane residential
street that sends a hundred vehicles an hour is not the same as a lane of
Garza Sada sending a thousand. A queue forms where demand exceeds capacity,
and it backs up from there. So we ask:

\begin{quote}
\emph{Where does the queue form on L\'azaro C\'ardenas, and is that place
predicted by the ratio of feeder lanes to corridor lanes, or by the
utilisation (demand over capacity) of each merge?}
\end{quote}

We expected utilisation to win, and we set a falsification rule before
running anything: if the lane ratio finds the simulated bottleneck at
least as often as utilisation does, the hypothesis fails.

The study makes four contributions. (i) An inventory of the corridor from
OpenStreetMap: every approach that can deliver vehicles, its lanes, class
and side, with internal express-to-frontage transfers removed
(Section~\ref{sec:data}). Its totals turn out to agree with the figures
that motivated the work. (ii) A queueing model at three levels: closed
forms validated by a vehicle-level simulator, a steady-state tandem
model, and a Cell Transmission Model with merges, partially FIFO
diverges and capacity drop (Section~\ref{sec:model}). (iii) A Monte
Carlo over every assumption we could not measure, scoring five
predictors of where the queue forms (Section~\ref{sec:results}).
(iv) Four interventions evaluated where the queue actually forms
(Section~\ref{sec:interventions}).
